\documentclass[a4paper]{scrreprt}
 
\usepackage[german]{babel}
\usepackage[utf8]{inputenc}
\usepackage[T1]{fontenc}
\usepackage{ae}
\usepackage[bookmarks,bookmarksnumbered]{hyperref}
\usepackage{graphicx}
\graphicspath{ {Images/} }
 
\begin{document}
 
    \begin{flushright}
        \includegraphics[scale = 0.7]{kit-logo.jpg}\\[0.5cm]
        % \includegraphics[scale = 1]{teco.jpg}
    \end{flushright}
    % \includegraphics[scale = 0.5]{kit-logo.jpg} \hspace{4cm} \includegraphics[scale = 1]{teco.jpg}
    \vspace*{2cm} 

    \begin{center} \large 
    
        Praxis der Softwareentwicklung
        \vspace * {1.5cm}

		{\bf \Huge Mind Rate}
		
		\vspace*{1cm}
		
        {\bf \large Ein interaktives Feedbacksystem f\"ur psychologische Untersuchungen nach Experience-Sampling-Method}
        \vspace*{1cm}

        {\large Pflichtenheft}
        \vspace*{2cm}

        Shanshan Du, Yi Ge, Renhan Lou, Ruoheng Ma, Haobin Tan
        \vspace*{1cm}

        02. Dezember 2016
        \vspace*{2.5cm}


        Betreuung: Anja Exler, Erik Psescara\\[1cm]
        Technology f\"ur Pervasive Computing\\[0.5cm]
        Karlsruher Institut für Technologie
    \end{center}
 
    % Die abstract-Umgebung kann bei Bedarf aus dem Pflichtenheft entfernt werden
    % \begin{abstract}
    %     Dies ist ein beispielhaftes Pflichtenheft in \LaTeX. Das Pflichtenheft
    %     beschreibt in konkreter Form, wie der Auftragnehmer die Anforderungen des
    %     Auftraggebers zu lösen gedenkt - das sogenannte wie und womit. Der Auftraggeber
    %     beschreibt vorher im Lastenheft möglichst präzise die Gesamtheit der
    %     Forderungen - was er entwickelt oder produziert haben möchte. Erst wenn der
    %     Auftraggeber das Pflichtenheft akzeptiert, sollte die eigentliche Arbeit beim
    %     Auftragnehmer beginnen.
    %     \vspace{1cm}
 
    %     Quelle: \url{http://de.wikipedia.org/wiki/Pflichtenheft} und Lehrbuch der
    %     Objektmodellierung von Heide Balzert
    %     \vspace{0.5cm}
 
    %     Quellcode: \url{http://www.karllorey.de/informatik-studium/vorlesungen/softwarepraktikum/pflichtenheft-in-latex/}
    % \end{abstract}
 
    % Platzierung des Inhaltsverzeichnisses
    \tableofcontents
 
    \chapter{Zielbestimmung}
    % Dieses Kapitel dient der Bestimmung von Zielen und nicht für deren Verwendung
    % notwendige Funktionen.
        Die Firma soll durch das Produkt in die Lage versetzt werden, psychologische Untersuchungen nach Experience-Sampling-Method(ESM) durchzuf\"uhren und ein Feedback der Untersuchungen zu erhalten.
 
 
        \section{Musskriterien}
            % Musskriterien: Für das Produkt unabdingbare Leistungen, die in jedem Fall
            % erfüllt werden müssen \footnote{gezwungen sein, etwas zu tun (Dies ist eine
            % beispielhafte Fußnote).}. Das System ist ohne diese Funktionen für seinen
            % gedachten Zweck nicht einsetzbar.

            \subsection{Verwalten von Teilnehmer der Untersuchung}
                \begin{itemize}
                    \item Teilnehmersstatus erfassen (z.B. Name, Email, Alter, Beruf)
                    \item Anmelden, Einloggen f\"ur Beteiligung der Untersuchung
                \end{itemize}
            
            \subsection{Verwalten der Frageb\"ogen}
                \begin{itemize}
                    \item Erstellen, \"Andern von Untersuchungsfragen
                    \item Erstellen, \"Andern von Scale und Optionen der Antworten
                    \item Setzen, \"Andern von G\"ultigkeitszeitbereich der Untersuchungsfragen
                    \item Setzen, \"Andern von H\"aufigkeit der Notifikationen zur Antworten
                \end{itemize}
            
            \subsection{Antwort auf Untersuchungsfragen}
                \begin{itemize}
                    \item regelm\"assige Notifikation zur Antworten
                    \item Untersuchungsfragen beantworten
                \end{itemize}
                
            \subsection{Verwalten von Antworten}
                \begin{itemize}
                    \item lokales Speichern von Antworten
                    \item Protokollieren von Zeitpunkt der Antwort
                    \item Hochladen von Antworten auf den Server
                \end{itemize}
            
            \subsection{Erfassen von Ergebnisse der Untersuchung}
                \begin{itemize}
                    \item Exportieren von Daten  
                \end{itemize}
 
        \section{Kannkriterien}
            Kannkriterien: Die Erfüllung der Kannkriterien ist erwünscht, jedoch nicht
            unbedingt notwendig. Sie sollten nur angestrebt werden, falls noch ausreichend
            Kapazitäten vorhanden sind.
 
        \section{Abgrenzungskriterien}
            Abgrenzungskriterien: Diese Kriterien sollen bewusst nicht erreicht werden.
 
    \chapter{Produkteinsatz}
        Der geplante Einsatz des Systems ist die Grundlage für Benutzungsoberfläche und
        Qualitätsanforderungen.
 
        \section{Anwendungsbereiche}
            Ein Pflichtenheft wird bspw. in einer IT-Abteilung genutzt.
 
        \section{Zielgruppen}
            Die Zielgruppe besteht also aus Informatikern, die mit der Projektplanung
            beauftragt wurden.
 
        \section{Betriebsbedingungen}
            Betriebsbedingungen: Die Betriebsbedingungen spezifiziert die physikalische
            Umgebung des Systems, die tägliche Betriebszeit, und ob das System ständiger
            Beobachtung durch Bediener ausgesetzt ist, oder ein unbeaufsichtigter Betrieb
            beabsichtigt ist.
 
    \chapter{Produktumgebung}
 
        \section{Software}
            Software: Gibt an, welche Software zum Betrieb vorhanden sein muss. Eine
            Aufteilung für Server und Client ist ggf. sinnvoll. Weiterhin sind unbedingt die
            kleinsten benötigten Versionsnummern anzugeben.
 
        \section{Hardware}
            Hardware: Hardware-Anforderungen des Systems.
 
    \chapter{Funktionale Anforderungen}
        Funktionalität: Spezifikation der einzelnen Produktfunktionen mit genauer und
        detaillierter Beschreibung.
 
        \begin{itemize}
            \item Typische Arbeitsabläufe
            \item Keine Angabe von typischen Verwaltungsfunktionen (CRUD \footnote{Create,
                  Read, Update, Delete}
        \end{itemize}
 
    \chapter{Produktdaten}
        Daten: Angabe der Daten, die langfristig aus Benutzersicht zu speichern sind.
 
    \chapter{Nichtfunktionale Anforderungen}
        Leistungen: Anforderungen bezüglich Zeit und Genauigkeit
 
    \chapter{Benutzungsoberfläche}
        Benutzungsoberfläche: grundlegende Anforderungen, Zugriffsrechte
 
        \begin{figure}[ht]
            \centering
            \rule{8cm}{6cm}
            \caption{Dies könnte ein Bild der Benutzungsoberfläche sein}
        \end{figure}

    \chapter{Globale Testfälle}
 
    \chapter{Qualitätsziele}
        Qualiätsziele: Allgemeine Ziele sind meistens Änderbarkeit und Wartbarkeit.
        Ziele sollten jedoch grundsätzlich messbar, spezifisch und relevant sein.
 
    \chapter{Glossar}
 
    % Abbildungsverzeichnis
    \listoffigures
 
\end{document}
